\scanpage{25}
\diagram{\begin{tikzcd}[ampersand replacement=\&,column sep=large]
\deleted{\mathrm{Hom}(\widetilde A,\widetilde X)}\arrow[r,"\deleted{u-\widetilde{\unsure{g|\eta}}\widetilde u}"]\arrow[ddr,"\deleted{\rho_A}"] \& \deleted{\mathrm{Hom}(\widetilde A,\widetilde Y)}\\
\deleted{\mathrm{Hom}(A,T(X))}\arrow[u,"\deleted{\sim}"]\arrow[r,"\deleted{u\to T(g)u}"] \& \deleted{\mathrm{Hom}(A,T(Y))}\arrow[u,"\deleted{\sim}"]\\
\&\deleted{\mathrm{Hom}(A,Y)}\arrow[u]
\end{tikzcd}}

\deleted{où $\beta_A$ est évidemment unique. Il revient au $=$ de dire que l'application \deleted{\illegible{1}} $T(g):T(X)\unsure{-|\to}T(Y)$ se factorise en $T(X)\longrightarrow Y\xrightarrow{i_Y}T(Y)$. \unsure{Disons|Disons nous|On voit}, que $g$ \uline{efface} le foncteur $T(X)/i_X(X)$ \deleted{\unsure{On convient que|Démonstration, que}}}

\deleted{\uline{Exemple} Si $g$ est un morphisme \unsure{spécial|quelconque} il en est de $=$ de $\lambda g$ quel que soit $\lambda\in\mathrm{Hom}(Y,\unsure{Z})$.}

\deleted{Un morphisme \deleted{\illegible{1}} \deleted{\unsure{$g$}} $i_M:M-T(M)$ est tjrs effaçant}

On appelle $t$-\deleted{résolution d'un} \inserted{complexe droit sur un} objet $A\in\mathcal C$ un complexe \inserted{droit} $X$ sur $A$, $A\to X$, muni de morphismes $t:T(A)-X^0$, $T(X^i)-X^{i+1}$ $(i>0)$ rendant commutatifs les diagrammes

\diagram{\begin{tikzcd}[ampersand replacement=\&]
A\arrow[r,"i_A"]\arrow[dr,"\varepsilon"'] \& T(A)\arrow[d,"t"]\\
\&X^0
\end{tikzcd}}

\diagram{\begin{tikzcd}[ampersand replacement=\&]
X^i\arrow[r,"i_{X^i}"]\arrow[dr,"d^i"'] \& T(X^i)\arrow[d,"t"]\\
\&X^{i+1}
\end{tikzcd}}

\deleted{Soit alors $t$} Soit donc $\unsure{A|A'}\to X$ une $s$-résolution \deleted{de} $A$, $B\to Y$ un $t$-complexe droit sur $B$, et $u:A\to B$

% paper watermark: ARGENTINE TYPEWRITER (vertical, ends cropped)
% archive sheet number: 24
